\chapter{State of the Art and Requirements Analysis}

This chapter is organized in two complementary parts. The first part reviews the current state of the art in digital agricultural advisory systems, identifying the main technological trends and the gaps that motivate this work. The second part translates these findings into a structured requirements specification for the AgriConnect platform, defining the stakeholders, use cases, and functional and non-functional constraints that guide its design.

\section{State of the Art}

\subsection{Digital Transformation of Agriculture in Morocco}

Agriculture remains a strategic sector in Morocco, yet it is increasingly constrained by structural challenges such as water scarcity, climate variability, fragmented farming systems, and unequal access to technical advisory services. These constraints make agricultural modernization not only a question of productivity, but also of resilience and resource efficiency. Morocco's \textit{Génération Green 2020--2030} strategy explicitly emphasizes the modernization of agriculture and the strengthening of human capital \cite{generationgreen2020}. In parallel, institutional analyses underline the importance of climate-resilient irrigation, improved water governance, and more accessible advisory support for farmers in the Moroccan context \cite{worldbank2024irrigation}.

At the same time, Morocco is undergoing a broader digital transformation that create    \7s favorable conditions for technology-enabled agricultural services. Recent national statistics show that mobile equipment is nearly universal among Moroccan households and that Internet access has expanded significantly in rural areas \cite{anrt2025}. This evolution makes mobile-based agricultural service delivery increasingly feasible, especially in rural environments where smartphones are more accessible than desktop-based platforms.

\subsection{Mobile-First Agricultural Advisory Systems}

Recent developments in digital agriculture show a transition from static information portals and conventional support channels toward mobile-first advisory systems. Such systems aim to provide practical and timely guidance directly through devices already embedded in farmers' daily routines. In developing agricultural contexts, accessibility is often as important as technical sophistication.

In Morocco, this mobile-first paradigm is particularly relevant because rural access to digital services relies predominantly on mobile connectivity. Messaging-based interaction therefore represents a more realistic interface than specialized web portals for many users. It lowers the adoption barrier and better aligns with common usage patterns, especially for farmers who may prefer sending short messages, voice notes, or crop images rather than navigating formal information systems \cite{anrt2025}.

\subsection{Conversational AI and Large Language Models in Agriculture}

Agricultural advisory systems initially relied on rule-based chatbots and predefined question--answer structures. While these systems were useful for handling simple and repetitive requests, they lacked flexibility when faced with informal language, ambiguous descriptions, multilingual communication, or multi-turn interactions. The rise of Large Language Models (LLMs) has substantially changed this landscape by enabling more natural dialogue, stronger language understanding, contextual reasoning, and adaptive response generation \cite{li2025llm}.

However, the use of LLMs in agriculture raises important concerns regarding reliability, factual accuracy, and domain specificity. Agricultural advice often involves high-stakes decisions related to irrigation, disease symptoms, pest management, or treatment timing. Consequently, the current state of the art favors LLM-based agricultural systems in which generative models are constrained by external knowledge sources, structured workflows, and expert supervision \cite{li2025llm,yin2025foundation}.

This requirement is particularly relevant in Morocco, where farmer interactions may involve multiple languages or mixed linguistic expressions, including Arabic, Darija, and French. Therefore, conversational agricultural systems intended for Moroccan users should not only be linguistically flexible, but also capable of producing grounded, context-aware, and cautious responses.

\subsection{Computer Vision for Plant and Crop Health Monitoring}

Computer vision has become a major field of research in smart agriculture, particularly for plant disease detection, pest recognition, and crop identification from field images. Advances in deep learning have enabled the development of image-based systems that support crop monitoring in a rapid and scalable manner. In practice, this is highly relevant because farmers frequently respond to visible crop problems by taking photographs with their smartphones \cite{yin2025foundation,zhu2025vlm}.

Despite these advances, the literature also highlights important limitations. Many image classification models are trained on controlled datasets that do not fully reflect real field conditions. Their performance may degrade in the presence of poor lighting, cluttered backgrounds, low image quality, or symptom overlap. Moreover, visible symptoms may be caused not only by plant diseases, but also by drought stress, salinity, nutrient deficiencies, or irrigation problems \cite{yin2025foundation,zhu2025vlm}.

These limitations are particularly relevant in the Moroccan context, where climatic stress and water scarcity can generate crop symptoms that visually resemble pathological disorders. Therefore, image-based analysis alone is insufficient for robust agricultural decision support, and current research increasingly treats visual analysis as one component within a broader advisory framework.

\subsection{Multimodal Artificial Intelligence in Agriculture}

To overcome the limitations of purely textual or purely visual systems, recent research increasingly explores multimodal artificial intelligence in agriculture. Multimodal systems combine several forms of input, such as text, images, and contextual signals, in order to improve reasoning quality and support more reliable outputs \cite{zhu2025vlm,yin2025foundation}.

In the Moroccan agricultural context, multimodal interaction is particularly relevant because farmers are more likely to provide an informal text message, a voice note, or a crop image than a structured technical report. A multimodal advisory system is therefore better aligned with actual usage patterns and can more effectively bridge the gap between unstructured farmer observations and actionable agronomic guidance.

\subsection{Retrieval-Augmented Generation for Trustworthy Agricultural Advice}

One of the main limitations of foundation models in specialized domains is their tendency to generate plausible but unsupported statements. In agriculture, such hallucinations are especially problematic because inaccurate recommendations may directly affect crop health, farm productivity, and input costs. This limitation has motivated the growing adoption of Retrieval-Augmented Generation (RAG), where generated responses are grounded in an external knowledge base \cite{li2025llm,yin2025foundation}.

RAG-based systems improve the trustworthiness of agricultural advisory platforms by linking responses to curated technical documents, agronomic guides, and domain-specific references. In the Moroccan context, this is especially important because useful agricultural advice often depends on region-specific practices, local crop conditions, and environmental constraints. A knowledge-grounded system is therefore more appropriate than a purely generative one, as it can provide more traceable and credible support.

\subsection{Climate- and Weather-Aware Agricultural Decision Support}

Agricultural recommendations cannot be fully reliable without accounting for environmental context. Irrigation planning, disease risk, spraying conditions, and crop stress management all depend strongly on weather conditions and seasonal variability. This issue is particularly significant in Morocco, where agriculture is highly exposed to drought, water scarcity, and climatic instability \cite{worldbank2024irrigation}.

For this reason, the integration of weather and climate information into digital advisory systems has become a major requirement in state-of-the-art agricultural decision support. A recommendation that ignores temperature peaks, rainfall probability, humidity conditions, or local stress indicators may have limited practical value. Consequently, modern agricultural support systems increasingly incorporate weather-aware reasoning in order to transform general agronomic knowledge into operational and localized advice \cite{worldbank2024irrigation}.

\subsection{Human-in-the-Loop Advisory and Expert Escalation}

Despite rapid advances in artificial intelligence, agriculture remains a high-stakes domain where full automation is rarely appropriate. Current research therefore emphasizes \textit{human-in-the-loop} approaches, in which AI systems assist with triage, information retrieval, and preliminary interpretation, while uncertain or critical cases are escalated to human experts \cite{li2025llm,yin2025foundation}.

This principle is particularly important in the Moroccan context, where trust in digital advisory tools depends strongly on their credibility and practical usefulness. Farmers are more likely to adopt such systems if they are capable of recognizing uncertainty and referring complex situations to agronomists instead of producing overconfident recommendations. Thus, hybrid architectures that combine intelligent automation with expert validation are among the most relevant models for operational agricultural support in Morocco.

\subsection{Positioning of the Project}

The literature demonstrates significant progress in mobile agricultural advisory, conversational AI, crop image analysis, multimodal reasoning, retrieval-based grounding, and weather-aware decision support. However, many existing solutions remain partial, focusing only on one dimension of the problem. Some systems specialize in plant image classification, others in generic conversational assistance, and others in document retrieval without personalization or escalation \cite{li2025llm,zhu2025vlm,yin2025foundation}.

The main gap lies in the lack of integrated systems capable of combining accessibility, contextual understanding, grounded knowledge, environmental awareness, and expert referral within a single workflow. In the Moroccan context, such integration is particularly relevant because it addresses the concrete realities of mobile usage, linguistic diversity, climate stress, and unequal access to technical expertise. From this perspective, the proposed project is aligned with the current state of the art by aiming to deliver a mobile-first, multilingual, multimodal, knowledge-grounded, weather-aware, and human-supervised agricultural support platform tailored to Moroccan farmers.

\subsection{Conclusion}

In summary, the state of the art in agricultural advisory is evolving toward integrated intelligent systems that combine conversational AI, computer vision, retrieval-based grounding, and environmental context. In Morocco, this evolution is especially relevant because agricultural decision-making is shaped by water scarcity, climate stress, regional diversity, and unequal access to technical support \cite{generationgreen2020,worldbank2024irrigation,anrt2025}. The increasing availability of mobile Internet in rural areas makes mobile-first advisory platforms particularly promising, while the limitations of purely generative or purely visual approaches justify the adoption of multimodal, knowledge-grounded, and human-supervised architectures \cite{anrt2025,li2025llm,zhu2025vlm,yin2025foundation}.

From this perspective, the development of an AI-assisted agricultural support platform tailored to Moroccan farmers is consistent with both international technological trends and national agricultural priorities.

\section{Problem Statement and Needs Analysis}

Agriculture in Morocco remains a strategic sector, yet many farmers continue to face practical barriers in obtaining timely and reliable agronomic support. These barriers include the late identification of crop diseases, the difficulty of accessing expert advice rapidly, the uneven availability of trustworthy scientific information, and the challenge of making correct decisions under uncertain environmental and climatic conditions. The project description explicitly identifies several of these difficulties, including crop diseases, water stress, unpredictable climatic conditions, limited rapid access to agronomic experts, and weak dissemination of scientific information among small-scale farmers \cite{generationgreen2020,worldbank2024irrigation}.

In real farming situations, the producer may observe unusual symptoms on leaves, fruits, or stems, but may not know whether the issue is urgent, whether it corresponds to a disease, a pest, or environmental stress, or whether immediate professional intervention is required. This uncertainty can lead either to delayed action or to inappropriate intervention. In this context, there is a clear need for an assistance system capable of receiving farmer input in a simple manner, interpreting it intelligently, and returning a response that is both understandable and operationally useful.

AgriConnect addresses this need by proposing a hybrid solution that combines accessibility and intelligence. On the one hand, WhatsApp serves as the main interaction channel, which reduces the technical barrier for farmers and makes the service more natural to use. On the other hand, the platform enriches this channel with intelligent processing modules able to structure farmer requests, analyze visual content, consult agricultural documentation, consider weather conditions, and escalate complex cases to human agronomists. The project description further specifies that farmers can consult their cases, agronomists can manage urgent interventions, and administrators can supervise the overall functioning of the platform \cite{anrt2025}.

Therefore, the need addressed by AgriConnect is broader than automatic answering. The system must create a coherent digital assistance environment in which artificial intelligence and human expertise complement each other. The workflow logic confirms this orientation, since incoming messages are normalized, buffered, enriched with context, and then routed toward direct response, knowledge retrieval, weather support, or escalation depending on the situation. 

\section{Stakeholders and Main Actors}

The AgriConnect platform involves three principal categories of stakeholders: farmers, agronomists, and administrators. Each of these actors has distinct objectives and interacts with the system in a different manner.

\subsection{Farmers}

Farmers are the primary users of the platform. Their main objective is to obtain quick, clear, and practically useful support regarding crop-related problems. They interact mainly through WhatsApp, where they can send text messages, voice messages, and images. The incoming message reception workflow confirms that the system is designed to handle these three main content types, normalize them, and prepare them for downstream reasoning \cite{anrt2025,zhu2025vlm}.

From the farmer's perspective, the platform must reduce uncertainty and provide accessible assistance. The farmer should be able to describe a field problem naturally, send a photo if needed, receive a preliminary response, and continue the interaction in a simple and intuitive way. In addition, the project description states that farmers should have access to a web interface to consult their submitted cases, check their resolution status, identify available agronomists, and provide feedback after a case is resolved \cite{li2025llm}.

\subsection{Agronomists}

Agronomists constitute the expert human layer of the system. Their role is to review urgent or complex situations that cannot be handled with sufficient confidence by the automated system. According to the project description, agronomists must be able to visualize urgent cases on an interactive map, take charge of cases according to their specialty and region, communicate with farmers when necessary, and access intervention history \cite{li2025llm,yin2025foundation}.

In practical terms, agronomists require a structured work environment rather than raw, unorganized user messages. The system must therefore provide them with case summaries, attached images, priority information, and relevant contextual elements, so that they can intervene more efficiently. This logic is also reflected in the escalation workflow, where a formal case is created and images may be attached when the situation requires human intervention \cite{li2025llm,yin2025foundation}.

\subsection{Administrators}

Administrators are responsible for the governance and supervision of the platform. The project description explicitly assigns them responsibilities such as user management, system supervision, statistical monitoring, and knowledge base management \cite{li2025llm}.

In addition to these general responsibilities, an especially important use case concerns the registration of agronomists. Unlike a standard user registration process, the onboarding of agronomists must be controlled because these users are expected to intervene as trusted agricultural experts. For this reason, when an agronomist registers on the platform, the administrator must review the submitted supporting documents, verify the authenticity and relevance of the profile, and then approve or reject the registration request. This administrative verification step is essential because it protects the professional reliability of the platform and helps ensure that only validated agronomists can assist farmers.

\section{Use Case Overview}

The main use cases of AgriConnect involve the interaction of the three actors identified above: the farmer, the agronomist, and the administrator. These use cases capture the essential services expected from the system and provide a functional overview of the platform.

The global use case diagram is presented in the system design chapter, where the system is modeled in full UML detail. The most important farmer-related use cases include sending a text message, sending an image, sending an audio message, receiving a preliminary automated reply, consulting case status, and giving feedback after resolution. On the agronomist side, the central use cases include registering on the platform, accessing urgent or assigned cases, reviewing farmer information, and managing interventions. On the administrator side, the most important use cases include managing users, supervising the platform, maintaining the knowledge base, and validating agronomist registration requests after document verification.

\section{Functional Requirements}

Table~\ref{tab:functional-req} summarizes the functional requirements grouped by actor.

\begin{table}[H]
\centering
\caption{Functional requirements by actor}
\label{tab:functional-req}
\begin{tabular}{p{2.5cm} p{3.5cm} p{7cm}}
\toprule
\textbf{Actor} & \textbf{Requirement} & \textbf{Description} \\
\midrule
\multirow{5}{*}{Farmer}
  & Multimodal input      & Send text, audio, and image messages via WhatsApp \cite{zhu2025vlm,yin2025foundation} \\
  & Automated reply       & Receive a contextualized response grounded in profile and history \cite{li2025llm} \\
  & Knowledge retrieval   & Get responses grounded in the agricultural knowledge base \cite{li2025llm,pgvector2024} \\
  & Weather context       & Benefit from weather-aware recommendations \cite{worldbank2024irrigation} \\
  & Case tracking         & Consult case status and provide feedback via the web app \cite{li2025llm} \\
\midrule
\multirow{3}{*}{Agronomist}
  & Registration          & Submit registration with documents; account remains pending until admin approval \\
  & Case management       & Access, review, and handle assigned cases with attached images and summaries \cite{li2025llm,yin2025foundation} \\
  & Intervention history  & Consult past interventions and update case status \\
\midrule
\multirow{3}{*}{Administrator}
  & User management       & Manage all user accounts and deactivate profiles when necessary \cite{li2025llm} \\
  & Agronomist validation & Review registration documents and approve or reject agronomist applications \\
  & KB maintenance        & Supervise platform activity and manage the knowledge base \\
\midrule
\multirow{4}{*}{System}
  & Message buffering     & Receive, normalize, and buffer WhatsApp interactions by conversation \\
  & Visual analysis       & Analyze crop images via the Vision Sub-Workflow and PlantNet API \\
  & RAG retrieval         & Chunk, embed, and retrieve knowledge base documents \cite{li2025llm,pgvector2024} \\
  & Escalation            & Create formal agronomist cases when automated resolution is insufficient \cite{li2025llm,yin2025foundation} \\
\bottomrule
\end{tabular}
\end{table}

\section{Non-Functional Requirements}

Table~\ref{tab:nonfunctional-req} presents the non-functional requirements and how each is addressed by the system design.

\begin{table}[H]
\centering
\caption{Non-functional requirements}
\label{tab:nonfunctional-req}
\begin{tabular}{p{3cm} p{5cm} p{5cm}}
\toprule
\textbf{Requirement} & \textbf{Description} & \textbf{How satisfied} \\
\midrule
Accessibility    & Usable by farmers with limited digital literacy & WhatsApp as main channel; no app installation required \cite{anrt2025} \\
Usability        & Concise, role-adapted responses and interfaces & LLM-generated replies adapted to farmer profile; role-specific web dashboards \\
Modularity       & Independent evolution of each service & Agent-based architecture; each workflow is independently deployable \\
Reliability      & Cautious behavior under uncertainty & Fallback preanalysis defaults; RAG grounding; escalation on low confidence \\
Traceability     & Persistent records of all interactions & Normalized batches, media files, and KB updates persisted in Supabase \\
Security         & Controlled access to data and expert roles & JWT-based auth; role-based access control; admin validation of agronomists \\
\bottomrule
\end{tabular}
\end{table}

\section{Chapter Summary}

This chapter clarified the needs addressed by AgriConnect and translated the project vision into a structured requirements specification. The analysis showed that the platform must support multimodal farmer interaction, provide contextual and grounded assistance, connect automated processing with expert human intervention, and offer dedicated interfaces for agronomists and administrators. It also highlighted a particularly important administrative use case: the verification of agronomist registration requests through document review before validation or rejection. These requirements now provide a clear reference framework for the design chapter, in which the architecture and internal organization of the system will be presented.